
\uline{plat} @@M0@@ \uline{tels que toute composante irréductible} @@M1@@ \uline{de} @@M2@@\uline{, munie de la structure réduite induite, soit intègre sur} @@M3@@.

En effet il existe une extension finie @@M4@@ du corps des fonctions @@M5@@ telle que toute composante irréductible de @@M6@@ soit géométriquement intègre (d'après EGA IV 4.6.8). Soit @@M7@@ un @@M8@@-schéma intègre tel que @@M9@@ et qu'il existe un ouvert non vide @@M10@@ de @@M11@@ et un morphisme fini plat @@M12@@. Il existe un ouvert non vide @@M13@@ de @@M14@@ tel que toute composante irréductible de @@M15@@, munie de la structure réduite induite, soit intègre sur @@M16@@ (EGA IV 9.7.7). Prenons @@M17@@ et @@M18@@.

Démontrons maintenant 1.2. Soient @@M19@@ et @@M20@@ comme dans 1.3. Pour tout @@M21@@, soit @@M22@@ une présentation de Chow avec @@M23@@ intègre et @@M24@@ birationnel. Quitte à restreindre @@M25@@ et @@M26@@, on peut supposer @@M27@@ intègre sur @@M28@@ (EGA IV 9.7.7). Soit @@M29@@ le morphisme canonique. Alors @@M30@@ est représentable par un schéma en groupes somme d'ouverts quasi-projectifs de type fini.

Enfin, 1.1 dit que quitte à restreindre @@M31@@ et @@M32@@, le morphisme @@M33@@ est représentable par un morphisme quasi-affine. Par suite @@M34@@ est représentable par un schéma en groupes, somme d'ouverts quasi-projectifs de type fini~; grâce à la descente plate finie (SGA 1 VIII 7.6 et EGA II 6.6.1), @@M35@@ l'est aussi.

\textbf{Lemme 1.4.} \uline{Soient} @@M36@@ \uline{un corps,} @@M37@@ \uline{et} @@M38@@ \uline{deux} @@M39@@\uline{-schémas en groupes localement de type fini et soit} @@M40@@ \uline{un morphisme quasi-affine. Alors} @@M41@@ \uline{est affine.}

En effet soit @@M42@@~; @@M43@@ est un @@M44@@-schéma en groupes de type fini qui est quasi-affine~; en vertu de (SGA 3 @@M45@@ 11.11), @@M46@@ est alors affine.
